\section{Method}\label{sec:method}

\subsection{Pipeline Overview}

Figure~\ref{fig:framework_overview} summarizes the intended pipeline. This section follows the same high-level order: ordered flow records are converted into sliding windows, windows are scored by the temporal model, critic and feature-deviation evidence are fused into one anomaly score, an independently calibrated threshold converts the score into an alarm, and post-hoc evidence is produced for analyst inspection.

We design the detector as a calibrated window-level alarm pipeline. The input is a temporally ordered sequence of network-flow records. Consecutive records are grouped into fixed-length windows \(X\in\mathbb{R}^{W\times F}\), where \(W\) is the number of time steps in a window and \(F\) is the number of flow features. The detector assigns one anomaly score to each window and then converts that score into an alarm using a threshold calibrated from benign data.

The main model is a Temporal Convolutional Network Wasserstein Generative Adversarial Network with Gradient Penalty (TCN-WGAN-GP). The TCN backbone models local temporal structure inside each window. The Wasserstein Generative Adversarial Network with Gradient Penalty (WGAN-GP) objective trains a critic that separates benign training windows from generated windows. On top of this critic, we compute a fused Temporal Attention-Deviation (TAD) score that combines critic evidence with deviation from benign reference statistics. The final threshold is not chosen on the test set; it is calibrated from an independent benign split and then fixed for final evaluation.

This method and its data protocol are coupled. Scaler fitting, model training, benign-reference construction, threshold calibration, and final testing use disjoint data roles. This design prevents the alarm threshold or reference statistics from using test information and makes the observed benign false-positive rate on the test split an operational outcome rather than a tuned quantity.

For implementation clarity, the end-to-end computation follows six fixed steps:
\begin{enumerate}
    \item Construct sliding windows from ordered flow records using \((W,r)\).
    \item Fit preprocessing and train the TCN-WGAN-GP model on model-train benign windows only.
    \item Compute critic-side and feature-deviation raw scores for each window.
    \item Normalize and fuse the two score branches into \(S_{\mathrm{TAD}}(X)\).
    \item Calibrate threshold \(\tau\) from an independent benign split and keep it fixed.
    \item Generate alarms and produce post-hoc audit traces through attribution and masking checks.
\end{enumerate}

\subsection{Sliding Windows and Data Roles}

Let the ordered flow sequence be \(\mathbf{X}=\{x_t\}_{t=1}^{T}\), where each \(x_t\in\mathbb{R}^{F}\) is a normalized network-flow feature vector and \(T\) is the number of ordered flow records. Given window length \(W\) and stride \(r\), the \(i\)-th window is
\begin{align}
X_i &= [x_{s_i}, x_{s_i+1}, \ldots, x_{s_i+W-1}]
      \in \mathbb{R}^{W\times F}, \\
s_i &= 1+(i-1)r .
\end{align}
Let \(y_i\in\{0,1\}\) be the evaluation label for \(X_i\), where \(y_i=1\) denotes an anomalous window. A window is labeled anomalous for evaluation if its attack-record ratio exceeds the predefined anomaly-ratio threshold. Labels are used only for reporting metrics after the threshold has been fixed. During model construction, only benign windows are used for feature scaling, model training, and TAD reference statistics.

The pipeline uses five data roles:
\begin{itemize}
    \item Feature scaler fitting uses model\_train\_benign only.
    \item Model training uses model\_train\_benign only.
    \item TAD reference construction uses model\_train\_benign only.
    \item Threshold calibration uses independent\_calibration\_benign only.
    \item Test labels are used only after threshold fixation for reporting metrics.
\end{itemize}

These roles are part of the method, not only an evaluation convention. The same anomaly score can appear stronger or weaker depending on how its threshold is selected. By separating the benign calibration split from the final test split, the detector is evaluated as a deployable fixed-threshold alarm generator.

\subsection{Attentive TCN-WGAN-GP Model}

The scoring model is the attentive TCN-WGAN-GP module. Let \(G_\theta\) be the generator with parameters \(\theta\), \(D_\phi\) be the critic with parameters \(\phi\), and \(z\) be a latent noise vector sampled from a prior distribution. A generated window is \(\tilde{X}=G_\theta(z)\). Real benign windows are sampled from \(P_r\), and generated windows induce \(P_g\). For the gradient penalty, \(\hat{X}\) is sampled along straight-line interpolations between real and generated windows. The critic and generator objectives are:
\begin{align}
\mathcal{L}_D
&=\mathbb{E}_{\tilde{X}\sim P_g}[D_\phi(\tilde{X})]-\mathbb{E}_{X\sim P_r}[D_\phi(X)] \nonumber\\
&\quad +\lambda\mathbb{E}_{\hat{X}}\left(\|\nabla_{\hat{X}}D_\phi(\hat{X})\|_2-1\right)^2,\\
\mathcal{L}_G
&=-\mathbb{E}_{\tilde{X}\sim P_g}[D_\phi(\tilde{X})].
\end{align}
Here, \(P_r\) is the empirical distribution of model-train benign windows, \(P_g\) is the generator distribution, and \(\lambda\) controls the gradient-penalty strength. The critic learns a scalar scoring function over windows. Compared with vanilla Generative Adversarial Network (GAN) training \cite{goodfellow2014gan}, WGAN-GP improves training stability and provides critic-side evidence that can be reused for anomaly scoring \cite{arjovsky2017wgan,gulrajani2017wgangp}.

\begin{figure*}[t]
  \centering
  \includegraphics[width=0.96\textwidth]{figures/model_architecture.png}
  \Description{Attentive TCN-WGAN-GP architecture showing generator/critic flows, WGAN-GP objective, feature-deviation branch, fused score, and final decision.}
  \caption{Attentive TCN-WGAN-GP architecture with critic evidence, feature-deviation scoring, and fused decision path. The diagram is illustrative; formal definitions and equations are given in this section.}
  \label{fig:model_architecture}
\end{figure*}

The critic includes attention pooling over temporal hidden states. Let \(h_t\) be the critic hidden state at time step \(t\), \(w_a\) and \(b_a\) be learnable attention parameters, and \(q_t\) be the unnormalized attention score. The normalized attention weight is:
\begin{equation}
a_t=\frac{\exp(q_t)}{\sum_{j=1}^{W}\exp(q_j)}, \qquad q_t=w_a^\top h_t+b_a.
\end{equation}
Here, \(a_t\) measures the contribution of time step \(t\) to the critic aggregation, and the weights sum to one over the \(W\) time steps. The pooled critic representation is
\begin{equation}
e(X)=\sum_{t=1}^{W}a_t h_t,
\end{equation}
where \(e(X)\) is the window-level embedding reused by the feature-deviation branch. The critic output \(D_\phi(X)\) is computed from the same temporal representation and provides critic-side evidence for whether a window resembles the benign training distribution. Attention is used as an internal temporal aggregation mechanism. We do not treat attention alone as a faithful explanation.

\subsection{Fused TAD Score}

For each window, the method combines two normalized anomaly components. The critic component \(S_D(X)\) reflects critic-side abnormality evidence from \(D_\phi\). In implementation, this branch uses the discriminator unknownness score
\begin{equation}
r_D(X)=1-\sigma(D_\phi(X)),
\end{equation}
where \(\sigma(\cdot)\) is the sigmoid function and larger \(r_D(X)\) means the discriminator assigns lower benign-likeness to the window.

The feature-deviation component \(S_F(X)\) measures how far the window embedding deviates from benign reference statistics computed only from model-train benign windows. Let \(\mathcal{B}_{\mathrm{train}}\) be the model-train benign windows and let
\begin{equation}
\mu_B=\frac{1}{|\mathcal{B}_{\mathrm{train}}|}\sum_{X\in\mathcal{B}_{\mathrm{train}}} e(X)
\end{equation}
be the benign embedding mean. The reported CICIDS2017 configuration uses the L2 feature-deviation score
\begin{equation}
r_F(X)=\|e(X)-\mu_B\|_2 .
\end{equation}
To place the two components on comparable scales, both raw scores are min-max normalized using benign-reference ranges from \(\mathcal{B}_{\mathrm{train}}\):
\begin{align}
S_D(X)&=\mathrm{clip}_{[0,1]}
\left(\frac{r_D(X)-m_D}{M_D-m_D+\epsilon}\right),\\
S_F(X)&=\mathrm{clip}_{[0,1]}
\left(\frac{r_F(X)-m_F}{M_F-m_F+\epsilon}\right),
\end{align}
where \(m_D,M_D\) are the minimum and maximum benign-reference critic scores, \(m_F,M_F\) are the minimum and maximum benign-reference feature-deviation scores, and \(\epsilon\) prevents division by zero. The fused Temporal Attention-Deviation score is:
\begin{equation}
S_{\mathrm{TAD}}(X)=\alpha S_D(X)+(1-\alpha)S_F(X),
\end{equation}
where \(S_D(X)\) is the normalized critic score, \(S_F(X)\) is the normalized feature-deviation score, and \(\alpha\in[0,1]\) is the mixing coefficient. In the selected strict operating configuration, \(\alpha=0.24\). Larger \(S_{\mathrm{TAD}}(X)\) indicates stronger anomaly evidence.

The computation path is therefore explicit: a window \(X\) is encoded by the TCN critic, attention pooling produces \(e(X)\), the critic branch computes \(r_D(X)\), the feature branch computes \(r_F(X)\), both branches are normalized using benign-reference statistics, and the fused TAD score is passed to threshold calibration.

\subsection{Independent Calibration and Alarm Decision}

The detector raises an alarm by comparing the fused score with a fixed threshold. Let \(\tau\) be the calibrated threshold and \(\hat{y}(X)\) be the predicted alarm label:
\begin{equation}
\hat{y}(X)=\mathbb{I}\left[S_{\mathrm{TAD}}(X)\ge\tau\right],
\end{equation}
where \(\mathbb{I}[\cdot]\) is the indicator function. The threshold \(\tau\) is selected from independent benign calibration scores under a target calibration false-positive rate \(\rho\). Equivalently, if \(Q_{1-\rho}\) denotes the empirical \((1-\rho)\)-quantile of the benign calibration scores, then:
\begin{equation}
\tau = Q_{1-\rho}\left(\{S_{\mathrm{TAD}}(X): X\in\mathcal{C}_{\mathrm{benign}}\}\right),
\end{equation}
where \(\mathcal{C}_{\mathrm{benign}}\) is the independent benign calibration split. The target \(\rho\) is therefore a threshold-selection parameter on calibration data. It is not the deployed false-positive rate. After \(\tau\) is fixed, deployment-facing alarm noise is measured by the observed benign false-positive rate on unseen test windows.

\subsection{Configuration and Sensitivity Scope}

The main CICIDS2017 configuration uses window size \(W=128\), stride \(r=16\), anomaly-ratio threshold \(0.15\), independent benign calibration ratio \(0.2\), attention pooling in the discriminator, fused score mode, and \(\alpha=0.24\). These settings define the reported operating pipeline; they are not re-tuned on the final test labels.

Sensitivity is evaluated with existing protocol-compatible checks rather than additional post-hoc tuning. The score ablation compares critic-only, feature-deviation-only, and fused TAD scoring under the same calibration rule. The threshold sweep varies the calibration-side target FPR while holding the independent calibration protocol fixed. The Gaussian perturbation analysis keeps the clean calibrated threshold fixed and measures how ranking metrics and fixed-threshold alarm metrics respond to input shifts. These analyses characterize the operating boundary of the fixed-threshold detector; they are not presented as exhaustive hyperparameter optimization or strong robustness evidence.

\subsection{Audit Evidence}

To support alarm inspection, we compute Integrated Gradients (IG) attributions \cite{sundararajan2017axiomatic} for scored windows. IG assigns attribution mass to time-feature positions with respect to the fused TAD score. We then compare attribution-guided masking with random masking: if masking top-ranked IG positions reduces the anomaly score more than masking random positions, the highlighted regions provide useful post-hoc evidence for the alarm decision.

This audit component is intentionally modest. It provides evidence localization and a decision-level sanity check for analysts, but it is not a proof of complete explanation faithfulness. When attention weights are available, they are reported as auxiliary temporal evidence and interpreted separately from IG.
